% Generated from locale/tr/content/first-order-logic/models-theories/expressing-relations.tex; do not hand-edit.


\olfileid[tr]{fol}{mat}{exr}

\olsection{\article{structure} \printtoken{S}{structure} İçindeki
  Bağıntıları İfade Etmek}

\begin{explain}
!!{formula}s ifadelerinin başlıca kullanım alanlarından biri,
!!a{structure}~$\Struct M$ içindeki özellik ve bağıntıları~$\Struct
M$ yapısının dili~$\Lang L$ içindeki temel simgeler cinsinden ifade
etmektir. Bununla şunu kastediyoruz: $\Struct M$ yapısının !!{domain}
bir nesneler kümesidir. !!{constant}s, !!{function}s ve !!{predicate}s
$\Struct M$ içinde sırasıyla $\Domain M$ içindeki bazı nesneler,
$\Domain M$ üzerindeki fonksiyonlar ve $\Domain M$ üzerindeki
bağıntılar olarak yorumlanır. Örneğin $\Obj A^2_0$, $\Lang L$ dilinde
bulunuyorsa $\Struct M$ ona $R=\Assign{{\Obj A^2_0}}{M}$ bağıntısını
atar. O zaman $\Atom{\Obj A^2_0}{\Obj v_1,\Obj v_2}$ formülü tam da
bu bağıntıyı şu anlamda \emph{ifade eder}: Bir değişken ataması~$s$,
$\Obj v_1$ değişkenini $a\in\Domain M$ nesnesine ve $\Obj v_2$
değişkenini $b\in\Domain M$ nesnesine eşliyorsa
\[
Rab \text{\quad ancak ve ancak\quad}
\Sat{M}{\Atom{\Obj A^2_0}{\Obj v_1, \Obj v_2}}[s].
\]
Burada değişken atamalarını kullanmak zorunda olduğumuza dikkat edin:
``$Rab$ ancak ve ancak $\Sat{M}{\Atom{\Obj A^2_0}{a,b}}$''
diyemeyiz; çünkü $a$ ve $b$ dilimizin simgeleri değil,
$\Domain M$ kümesinin !!{element}s ifadeleridir.

Yalnızca atomik !!{formula}s ile sınırlı olmadığımız; bunları mantıksal
bağlaçlar ve niceleyicilerle birleştirebildiğimiz için, daha karmaşık
!!{formula}s doğrudan $\Struct M$ yapısına yerleştirilmemiş başka
bağıntıları tanımlayabilir. Bunun nasıl yapıldığıyla ve özellikle
!!a{structure} içinde hangi bağıntıları tanımlayabileceğimizle
ilgileniyoruz.
\end{explain}

\begin{defn}
$!A(\Obj v_1,\dots,\Obj v_n)$, $\Lang L$ dilinde yalnızca
$\Obj v_1,\dots,\Obj v_n$ değişkenlerinin serbest geçtiği !!a{formula}
ve $\Struct M$,~$\Lang L$ için !!a{structure} olsun.
$!A(\Obj v_1,\dots,\Obj v_n)$ formülü, herhangi bir $s(\Obj v_i)=a_i$
($i=1,\dots,n$) değişken ataması için
\[
Ra_1\dots a_n \text{\quad ancak ve ancak\quad}
\Sat{M}{\Atom{!A}{\Obj v_1,\dots, \Obj v_n}}[s]
\]
olduğunda ve ancak bu durumda $R\subseteq\Domain M^n$ bağıntısını
\emph{ifade eder}.
\end{defn}

\begin{ex}
Aritmetiğin standart modeli~$\Struct N$ içinde $\Obj v_1<\Obj v_2
\lor\eq[\Obj v_1][\Obj v_2]$ !!{formula}, $\Nat$ üzerindeki $\le$
bağıntısını ifade eder. $\eq[\Obj v_2][\Obj v_1']$ !!{formula}, ardıl
bağıntısını; yani $m$, $n$'nin ardılı olduğunda $Rnm$ bağıntısının
geçerli olduğu $R\subseteq\Nat^2$ bağıntısını ifade eder.
$\eq[\Obj v_1][\Obj v_2']$ formülü öncül bağıntısını ifade eder.
$\lexists[\Obj v_3][(\eq/[\Obj v_3][\Obj 0]\land
\eq[\Obj v_2][(\Obj v_1+\Obj v_3)])]$ ve
$\lexists[\Obj v_3][\eq[(\Obj v_1+{\Obj v_3}')][v_2]]$
!!{formula}s ifadelerinin ikisi de $\Obj <$ bağıntısını ifade eder.
Bu, $<$ yüklem simgesinin aritmetik dilinde aslında gereksiz olduğu;
tanımlanabileceği anlamına gelir.
\end{ex}

\begin{explain}
Bu düşünce yalnızca belirli !!{structure}s içinde değil, bir dili
amaçlanan bir modeli ya da modelleri betimlemek için kullandığımız;
yani kuramları ele aldığımız her durumda ilginçtir. Bu kuramlar temel
simgeler olarak çoğu kez yalnızca birkaç !!{predicate}s içerir; fakat
betimledikleri evrende başka birçok bağıntı önemli rol oynar. Bu başka
bağıntılar, dilin temel !!{predicate}s ifadelerini yorumlayan bağıntılar
aracılığıyla sistematik biçimde ifade edilebiliyorsa bunların dilde
\emph{tanımlanabildiğini} söyleriz.
\end{explain}

\begin{prob}
$\Lang L_A$ dilinde aşağıdaki bağıntıları tanımlayan !!{formula}s bulun:
\begin{enumerate}
\item $n$, $i$ ile $j$ arasındadır;
\item $n$, $m$'yi tam böler (yani $m$, $n$'nin bir katıdır);
\item $n$ bir asal sayıdır (yani $1$ ve $n$ dışında hiçbir sayı $n$'yi
  tam bölmez).
\end{enumerate}
\end{prob}

\begin{prob}
$!A(\Obj v_1,\Obj v_2)$ formülünün !!a{structure}~$\Struct M$ içinde
$R\subseteq\Domain M^2$ bağıntısını ifade ettiğini varsayın. Aşağıdaki
bağıntıları ifade eden formüller bulun:
\begin{enumerate}
\item $R$ bağıntısının tersi~$R^{-1}$;
\item bağıntı bileşkesi~$R\mid R$.
\end{enumerate}
$R$ bağıntısının geçişli kapanışı~$R^+$ ifadesini belirtmenin bir
yolunu bulabilir misiniz?
\end{prob}

\begin{prob}
$\Lang L$, yalnızca iki yerli $<$ yüklem simgesini içeren dil olsun
(elbette~$\eq$ dışında başka !!{constant}s, !!{function}s ya da
!!{predicate}s bulunmasın). $\Domain N=\Nat$ ve
$\Assign{<}{N}=\Setabs{\tuple{n,m}}{n<m}$ olan yapıya $\Struct N$
deyelim. Aşağıdakileri ispatlayın:
\begin{enumerate}
\item $\{0\}$, $\Struct N$ içinde tanımlanabilir;
\item $\{1\}$, $\Struct N$ içinde tanımlanabilir;
\item $\{2\}$, $\Struct N$ içinde tanımlanabilir;
\item her $n\in\Nat$ için $\{n\}$ kümesi $\Struct N$ içinde
  tanımlanabilir;
\item $\Domain N$ kümesinin her sonlu alt kümesi $\Struct N$ içinde
  tanımlanabilir;
\item $\Domain N$ kümesinin tümleyeni sonlu olan her alt kümesi
  $\Struct N$ içinde tanımlanabilir (burada $X\subseteq\Nat$ kümesinin
  tümleyeni sonludur ancak ve ancak $\Nat\setminus X$ sonludur).
\end{enumerate}
\end{prob}

